\documentclass[12pt,a4paper]{article}
\usepackage[utf8]{inputenc}
\usepackage[T1]{fontenc}
\usepackage[spanish]{babel}
\usepackage{amsmath,amssymb,amsfonts}
\usepackage{geometry}
\usepackage{booktabs}
\usepackage{hyperref}
\geometry{margin=2.5cm}

\title{Perturbaciones Tensoriales desde el Borde del Horizonte: Espectro de Ondas Gravitacionales y Relaciones de Consistencia}
\author{Kevin Mu\~noz}
\date{}

\begin{document}
\maketitle

\begin{abstract}
Los Pasos 1--4 del Marco del Horizonte Causal derivaron el espectro de potencias escalar primordial desde el horizonte aparente. El Paso~5 identificó la razón tensor-a-escalar $r$ como la predicción diferencial más aguda, condicional al Paso~6. Este documento cierra esa brecha. Mostramos que las dos polarizaciones transversas sin traza (TT) de la onda gravitacional, al restringirse a la dos-esfera del horizonte aparente, definen campos tensoriales de borde con variable canónica $\mu_k^\lambda = (M_{\rm Pl}/2)\,a(\eta)\,h_k^\lambda(\eta)$ --- la misma normalización que la función de modo tensorial del bulk. La ecuación de modo para $\mu_k^\lambda$ es la ecuación de Mukhanov--Sasaki con $z_T = a$, en contraste con el caso escalar donde $z_A = a/\sqrt\varepsilon$. Esta diferencia estructural tiene un origen físico transparente: el campo escalar de borde $\psi = \delta R_A/R_A$ es una perturbación de TAMAÑO del horizonte (cambio fraccional en el radio), que está suprimida por rodadura lenta porque el horizonte apenas se mueve durante la inflación; el campo tensorial $h^\lambda$ es una perturbación de FORMA (cizallamiento transverso) que no altera la tasa de expansión ni el radio del horizonte, y por tanto no lleva supresión de $\varepsilon$. El índice espectral tensorial resultante es $n_T = -2\varepsilon$, el espectro de potencias tensorial es $\Delta_h^2 = 2H^2/(\pi^2 M_{\rm Pl}^2)$, y la razón tensor-a-escalar es $r = 16\varepsilon$, en completo acuerdo con el resultado estándar de rodadura lenta. La relación de consistencia estándar $r = -8n_T$ se satisface. Para el CHF mínimo ($\beta_S = 0$), se obtiene una relación de consistencia adicional: $n_T = n_s - 1$, que implica $r = 8(1-n_s)$. Dado que $n_s = 0.9649$ (Planck 2018) da $r \approx 0.281 > 0.056$, el CHF mínimo queda excluido al 95\% CL por la cota actual $r < 0.056$. El marco requiere $\beta_S < -0.042$ para lograr compatibilidad simultánea con el $n_s$ observado y la cota sobre $r$.
\end{abstract}

\section{Motivación}

El Paso~5 identificó tres predicciones diferenciales del CHF. La más aguda fue la línea de consistencia $(n_s, r)$: $r = 8(1-n_s)$, que requirió asumir $r = 16\varepsilon$ --- un resultado tomado del cálculo estándar en el bulk y aún no derivado dentro del CHF. El Paso~5 dejó abiertas dos posibilidades: o bien el espectro tensorial del CHF coincide con el resultado estándar ($r = 16\varepsilon$ confirmado, marco mínimo falsificado), o bien difiere (relación de consistencia estándar rota, marco viable). El presente documento resuelve esto derivando el espectro de potencias tensorial directamente desde el borde del horizonte.

La estructura de la derivación es paralela al Paso~1 (cuantización escalar). Las dos polarizaciones tensoriales se identifican como campos en la dos-esfera del horizonte, se deriva una acción de borde por reducción dimensional desde la acción de Einstein--Hilbert en el bulk, y la ecuación de modo se obtiene mediante la identificación de modos comóviles $l_k(t) = k/(a(t)H(t))$ introducida en el Paso~1. El nuevo ingrediente clave es la distinción física entre perturbaciones escalares (de tamaño) y tensoriales (de forma) del horizonte.

\section{Grados de Libertad Tensoriales en el Horizonte Aparente}

\subsection{Identificando los modos tensoriales de borde}

En la descomposición ADM 3+1, las perturbaciones métricas se dividen en sectores escalar, vectorial y tensorial (TT). El sector tensorial está descrito por la parte transversa sin traza $h_{ij}^{TT}$ de la perturbación de la métrica espacial, que satisface $\partial^i h_{ij}^{TT} = 0$ y $h^{TT}_{ii} = 0$. Tiene dos grados de libertad físicos, correspondientes a los dos estados de polarización $\lambda = +, \times$.

El horizonte aparente $\mathcal{H}$ es una dos-esfera de radio $R_A(t) = 1/H(t)$ inmersa en la hipersuperficie espacial $\Sigma_t$. Una perturbación TT en el bulk $h_{ij}^{TT}$ induce una deformación métrica en $\mathcal{H}$:
\begin{equation}
\delta\gamma_{AB} = h_{ij}^{TT}\,e^i_A\,e^j_B \tag{1}
\end{equation}
donde $e^i_A$ son los vectores tangentes a la dos-esfera y $\gamma_{AB}$ es la métrica redonda. La deformación $\delta\gamma_{AB}$ es sin traza ($\gamma^{AB}\delta\gamma_{AB} = 0$) y transversa en $S^2$, representando una deformación de FORMA genuina de la dos-esfera.

\subsection{Identificación de modos}

Cada modo de polarización $h_k^\lambda(t)$ con número de onda comóvil $k$ se identifica con el modo angular $l_k(t) = k/(a(t)H(t))$ en la dos-esfera del horizonte, exactamente como en el caso escalar (Paso~1, ec.~(2)). El congelamiento ocurre cuando $l_k = 1$, es decir, $k = a(t_k)H$, que es la condición estándar de cruce del horizonte.

\section{La Acción Tensorial de Borde}

\subsection{Reducción dimensional desde el bulk}

La acción de Einstein--Hilbert en el bulk, expandida a segundo orden en $h_{ij}^{TT}$ alrededor del fondo FRW, da para cada modo de polarización $h_k^\lambda(\eta)$ en tiempo conforme:
\begin{equation}
S_T^{(k,\lambda)} = \frac{M_{\rm Pl}^2}{4}\int d\eta\, a^2\!\left[(h_k^{\lambda\prime})^2 - k^2(h_k^\lambda)^2\right] \tag{2}
\end{equation}
donde las primas denotan $d/d\eta$.

\subsection{Variable canónica}

Definiendo la variable canónica tensorial
\begin{equation}
\mu_k^\lambda \equiv \frac{M_{\rm Pl}}{2}\,a(\eta)\,h_k^\lambda(\eta) \tag{3}
\end{equation}
y sustituyendo en (2), la integración por partes da:
\begin{equation}
S_T^{(k,\lambda)} = \frac{1}{2}\int d\eta\!\left[(\mu_k^{\lambda\prime})^2 - \left(k^2 - \frac{a''}{a}\right)(\mu_k^\lambda)^2\right] \tag{4}
\end{equation}
La ecuación de movimiento es:
\begin{equation}
(\mu_k^\lambda)'' + \left[k^2 - \frac{a''}{a}\right]\mu_k^\lambda = 0 \tag{5}
\end{equation}
Esta es la \textbf{ecuación de Mukhanov--Sasaki para modos tensoriales}, con campo bomba $z_T = a$.

\subsection{Comparación con el caso escalar}

En el Paso~1, la variable canónica escalar era $u_k = z_A\psi_k$ con $z_A = a/\sqrt\varepsilon$. Para modos tensoriales, el campo bomba es simplemente $z_T = a$. Específicamente:
\begin{equation}
\frac{a''}{a} = \mathcal{H}^2(2 + \varepsilon + \cdots), \qquad \frac{z_A''}{z_A} = \mathcal{H}^2(2 + 3\varepsilon + \cdots) \tag{6}
\end{equation}
donde $\mathcal{H} = aH$. Esta diferencia a nivel subleader en rodadura lenta es el origen de los valores distintos de $\nu$ para modos escalares y tensoriales.

\section{Origen Físico de $z_T = a$ frente a $z_A = a/\sqrt\varepsilon$}

\paragraph{Modos escalares ($z_A = a/\sqrt\varepsilon$).} El campo escalar de borde es $\psi = \delta R_A/R_A$ --- una perturbación fraccional del RADIO del horizonte. Durante la rodadura lenta ($\varepsilon \ll 1$), el horizonte está casi congelado. Una perturbación de curvatura comóvil dada $\mathcal{R}$ induce solo un pequeño cambio fraccional $\psi \approx \varepsilon\mathcal{R}$ en el tamaño del horizonte (Paso~2). La supresión $\psi \ll \mathcal{R}$ está capturada por el $1/\sqrt\varepsilon$ en $z_A$.

\paragraph{Modos tensoriales ($z_T = a$).} El campo tensorial de borde $h^\lambda$ es una perturbación transversa sin traza de FORMA de la dos-esfera del horizonte --- deforma la esfera en un elipsoide sin cambiar el radio del horizonte. Dado que $h^\lambda$ satisface $h^{TT}_{ii} = 0$ y $\partial^i h_{ij}^{TT} = 0$, no lleva perturbación de la tasa de expansión local $H$. Una perturbación de forma no desencadena un desplazamiento en $R_A$ y por tanto no está suprimida por rodadura lenta. La normalización canónica está fijada enteramente por $M_{\rm Pl}$, sin factor $\varepsilon$.

\section{Vacío de Bunch--Davies y Congelamiento Tensorial}

\subsection{Potencial efectivo}

Para rodadura lenta con $\varepsilon = -\dot H/H^2 \ll 1$ y $\dot\varepsilon \ll H\varepsilon$:
\begin{equation}
\frac{a''}{a} = \frac{\nu_T^2 - \tfrac{1}{4}}{\eta^2}, \qquad \nu_T = \frac{3}{2} + \varepsilon + \mathcal{O}(\varepsilon^2) \tag{7}
\end{equation}
La ecuación de modo tensorial (5) se convierte en:
\begin{equation}
(\mu_k^\lambda)'' + \left[k^2 - \frac{\nu_T^2 - \tfrac{1}{4}}{\eta^2}\right]\mu_k^\lambda = 0 \tag{8}
\end{equation}
Esta tiene la misma forma que la ecuación escalar (Paso~1, ec.~(17)) con $\nu_S = 3/2 + \varepsilon - \beta_S/3$. Para $\beta_S = 0$: $\nu_S = \nu_T = 3/2 + \varepsilon$. Las dos ecuaciones de modo son idénticas a este orden.

\subsection{Condiciones iniciales de Bunch--Davies}

En tiempos tempranos ($|k\eta| \gg 1$), la ecuación (8) se reduce a un oscilador libre. La condición de Bunch--Davies selecciona el modo de frecuencia positiva:
\begin{equation}
\mu_k^\lambda(\eta) \xrightarrow{|k\eta|\to\infty} \frac{1}{\sqrt{2k}}\,e^{-ik\eta} \tag{9}
\end{equation}

\subsection{Solución exacta y amplitud en el congelamiento}

La solución exacta de (8) con condición (9) es:
\begin{equation}
\mu_k^\lambda(\eta) = \frac{\sqrt{\pi}}{2}\,e^{i(\nu_T + \tfrac{1}{2})\frac{\pi}{2}}\,\sqrt{-\eta}\,H_{\nu_T}^{(1)}(-k\eta) \tag{10}
\end{equation}
En escalas super-horizonte ($|k\eta| \to 0$):
\begin{equation}
|\mu_k^\lambda(\eta)|^2 \xrightarrow{|k\eta|\to 0} \frac{2^{2\nu_T - 3}\,\Gamma(\nu_T)^2}{\pi k}\left(\frac{1}{-k\eta}\right)^{2\nu_T - 1} \tag{11}
\end{equation}
El congelamiento ocurre en $k\eta_k = -1$. Para de Sitter ($\nu_T = 3/2$):
\begin{equation}
|\mu_k^\lambda|_{\rm congel}^2 = \frac{1}{2k} \tag{12}
\end{equation}
La amplitud tensorial en el congelamiento:
\begin{equation}
|h_k^\lambda|_{\rm congel} = \frac{2|\mu_k^\lambda|_{\rm congel}}{M_{\rm Pl}\,a_{\rm congel}} = \frac{2}{M_{\rm Pl}}\cdot\frac{H}{k}\cdot\frac{1}{\sqrt{2k}} = \frac{\sqrt{2}\,H}{M_{\rm Pl}\,k^{3/2}} \tag{13}
\end{equation}
usando $a_{\rm congel} = k/H$ en el cruce del horizonte.

\section{Espectro de Potencias Tensorial}

El espectro de potencias tensorial adimensional, sumado sobre ambas polarizaciones, es:
\begin{equation}
\Delta_h^2 \equiv \frac{k^3}{2\pi^2}\sum_\lambda |h_k^\lambda|_{\rm congel}^2 = \frac{k^3}{2\pi^2}\cdot\frac{2H^2}{M_{\rm Pl}^2\,k^3} \tag{14}
\end{equation}
\begin{equation}
\boxed{\Delta_h^2 = \frac{2H^2}{\pi^2 M_{\rm Pl}^2}} \tag{15}
\end{equation}
Este es exactamente el resultado estándar de inflación con rodadura lenta para la amplitud de ondas gravitacionales primordiales.

\paragraph{Índice espectral tensorial.}
De (11), la dependencia en $k$ del espectro tensorial a orden líder en rodadura lenta:
\begin{equation}
n_T \equiv \frac{d\ln\Delta_h^2}{d\ln k} = 3 - 2\nu_T = -2\varepsilon \tag{16}
\end{equation}

\section{Razón Tensor-a-Escalar}

Combinando el espectro tensorial (15) con el espectro de potencias escalar de la perturbación de curvatura comóvil $\mathcal{R}$:
\begin{equation}
\Delta_\mathcal{R}^2 = \frac{H^2}{8\pi^2\varepsilon M_{\rm Pl}^2} \tag{17}
\end{equation}
la razón tensor-a-escalar es:
\begin{equation}
r \equiv \frac{\Delta_h^2}{\Delta_\mathcal{R}^2} = \frac{2H^2/(\pi^2 M_{\rm Pl}^2)}{H^2/(8\pi^2\varepsilon M_{\rm Pl}^2)} \tag{18}
\end{equation}
\begin{equation}
\boxed{r = 16\varepsilon} \tag{19}
\end{equation}
Esta es la relación de consistencia estándar de inflación de campo único con rodadura lenta, derivada aquí directamente desde el borde del horizonte.

\section{Relaciones de Consistencia Estándar}

\subsection{La relación de Lyth--Liddle}

Combinando (16) y (19):
\begin{equation}
r = -8n_T \tag{20}
\end{equation}
Esta es la relación de consistencia estándar de inflación de campo único. El CHF la satisface porque tanto $r = 16\varepsilon$ como $n_T = -2\varepsilon$ se derivan del mismo parámetro termodinámico $\varepsilon$.

\subsection{Relación de consistencia adicional del CHF}

Para el CHF mínimo con $\beta_S = 0$ (escalar) y $\beta_T = 0$ (tensorial):
\begin{equation}
n_s - 1 = -2\varepsilon = n_T \tag{21}
\end{equation}
Los índices espectrales escalar y tensorial son iguales. Combinando con (20):
\begin{equation}
r = -8n_T = -8(n_s - 1) = 8(1 - n_s) \tag{22}
\end{equation}
Esto colapsa la libertad bidimensional de la inflación estándar a una única línea en el plano $(n_s, r)$.

\paragraph{Comparación con la inflación estándar.} En la inflación estándar de campo único con dos parámetros de rodadura lenta $\varepsilon$ y $\eta'$:
\begin{equation}
n_s - 1 = -2\varepsilon - \eta', \qquad n_T = -2\varepsilon \tag{23}
\end{equation}
por lo que $n_T - (n_s - 1) = \eta' \neq 0$ en general. El CHF predice $\eta' = 0$, lo cual es violado por la mayoría de los potenciales inflacionarios.

\section{Resolución de la Tensión del Paso~5}

El Paso~5 identificó una tensión: el CHF mínimo ($\beta_S = 0$) predice $r = 8(1-n_s) \approx 0.281$ para $n_s = 0.9649$, superando la cota de Planck 2018 $r < 0.056$. La presente derivación confirma que $r = 16\varepsilon$ sí se cumple. La tensión es real.

\subsection{Campo escalar no mínimo ($\beta_S \neq 0$)}

Para $\beta_S \neq 0$: $n_s - 1 = -2\varepsilon + 2\beta_S/3$, dando:
\begin{equation}
r = 16\varepsilon = 8(1 - n_s) + \frac{16\beta_S}{3} \tag{24}
\end{equation}
Para $n_s = 0.9649$ y $r < 0.056$:
\begin{equation}
\beta_S < -\frac{3}{16}(0.281 - 0.056) \approx -0.042 \tag{25}
\end{equation}
Se requiere un campo escalar de horizonte ligeramente taquiónico ($m_{\rm eff}^2 = \beta_S H^2 < 0$). Un $\beta_S$ negativo pequeño es físicamente plausible y no desestabiliza la ecuación de modo.

\subsection{Interpretación de la tensión como falsificación}

Si experimentos futuros establecen $r < 0.01$ (objetivo CMB-S4/LiteBIRD), esto requeriría $\beta_S < -0.065$. La fórmula del tilt da entonces:
\begin{equation}
n_s - 1 = -2\varepsilon + \frac{2\beta_S}{3} < -2\varepsilon - 0.043 \tag{26}
\end{equation}
Para $n_s = 0.965$: $\varepsilon > 0.004$, dando $r = 16\varepsilon > 0.064$ --- una contradicción con $r < 0.01$. El CHF mínimo queda excluido; el caso no mínimo permanece viable pero pierde simplicidad predictiva.

\section{Resumen de Resultados Espectrales}

\begin{table}[h]
\centering
\small
\begin{tabular}{lll}
\toprule
Cantidad & Resultado CHF & Inflación estándar \\
\midrule
Campo bomba tensorial & $z_T = a$ & $z = aM_{\rm Pl}$ (igual) \\
Índice espectral tensorial & $n_T = -2\varepsilon$ & $n_T = -2\varepsilon$ \\
Espectro de potencias tensorial & $\Delta_h^2 = 2H^2/(\pi^2 M_{\rm Pl}^2)$ & $\Delta_h^2 = 2H^2/(\pi^2 M_{\rm Pl}^2)$ \\
Relación de consistencia & $r = -8n_T$ & $r = -8n_T$ \\
Razón tensor-a-escalar & $r = 16\varepsilon$ & $r = 16\varepsilon$ \\
Relación adicional CHF (mínimo) & $n_T = n_s - 1$ & $n_T = n_s - 1 + \eta'$ ($\eta' \neq 0$ en general) \\
\bottomrule
\end{tabular}
\end{table}

\section{Conclusión}

Los grados de libertad tensoriales del Marco del Horizonte Causal son las dos polarizaciones TT de la onda gravitacional, tratadas como perturbaciones de forma de la dos-esfera del horizonte aparente. Su acción de borde se obtiene de la acción de Einstein--Hilbert en el bulk por reducción dimensional, con variable canónica $\mu_k^\lambda = (M_{\rm Pl}/2)a h_k^\lambda$. La ecuación de modo tiene campo bomba $z_T = a$, en contraste con el campo bomba escalar $z_A = a/\sqrt\varepsilon$. Con condiciones iniciales de Bunch--Davies, la ecuación de modo tensorial da $n_T = -2\varepsilon$, $\Delta_h^2 = 2H^2/(\pi^2 M_{\rm Pl}^2)$ y $r = 16\varepsilon$. Las relaciones de consistencia estándar se satisfacen. La relación específica del CHF $n_T = n_s - 1$ (para $\beta_S = 0$) confirma la predicción del Paso~5 $r = 8(1-n_s)$, colocando al CHF mínimo en tensión con la cota de Planck 2018 $r < 0.056$. La compatibilidad con los datos actuales requiere $\beta_S < -0.042$.

\begin{thebibliography}{9}
\bibitem{planck2018} Planck Collaboration, ``Planck 2018 results. X. Constraints on inflation,'' \textit{A\&A} \textbf{641}, A10 (2020).
\bibitem{bicep2021} BICEP/Keck Collaboration, ``BK18: Improved constraints on primordial gravitational waves,'' \textit{Phys. Rev. Lett.} \textbf{127}, 151301 (2021).
\bibitem{mukhanov1992} V. F. Mukhanov, H. A. Feldman y R. H. Brandenberger, ``Theory of cosmological perturbations,'' \textit{Phys. Rep.} \textbf{215}, 203 (1992).
\bibitem{lyth1997} D. H. Lyth, ``What would we learn by detecting a gravitational wave signal in the cosmic microwave background anisotropy?,'' \textit{Phys. Rev. Lett.} \textbf{78}, 1861 (1997).
\bibitem{munoz2025a} K. Mu\~noz, ``A Horizon-Centered Formalization of Cosmological Dynamics and Perturbations,'' Causal Horizon Framework, Documento~5 (2025).
\bibitem{munoz2025b} K. Mu\~noz, ``Scale Invariance from Horizon Field Quantization,'' Causal Horizon Framework (2025).
\bibitem{munoz2025c} K. Mu\~noz, ``Gauge-Invariant Formulation of the Horizon Perturbation,'' Causal Horizon Framework (2025).
\bibitem{munoz2025d} K. Mu\~noz, ``Spectral Tilt from Horizon Thermodynamics,'' Causal Horizon Framework (2025).
\bibitem{munoz2025e} K. Mu\~noz, ``Differential Predictions of the Causal Horizon Framework,'' Causal Horizon Framework (2025).
\end{thebibliography}

\end{document}
